\chapter{Notacion y supuestos comunes}
\label{ch:notacion}

\section{Variables de momento}

\begin{align}
\vk &:\ \text{momento de Bloch o variable global del espacio reciproco},\\
\vq &:\ \text{desviacion pequena respecto de un nodo o punto de expansion},\\
\vp &:\ \text{momento efectivo continuo en el sector infrarrojo}.
\end{align}

Regla practica:
\begin{enumerate}[label=\arabic*.]
  \item \vk\ cuando hablemos del modelo discreto exacto;
  \item $\vk = \vk_0 + \vq$ al linealizar alrededor de un nodo;
  \item \vp\ una vez pasado al lenguaje efectivo continuo.
\end{enumerate}

\section{Espacios internos}

\begin{align}
\sigma_i &:\ \text{grado de libertad tipo Weyl / pseudospin de dos componentes},\\
\tau_i   &:\ \text{indice quiral o bloque Weyl acoplado}.
\end{align}

La forma canonica del bloque Dirac efectivo queda:
\begin{equation}
\HD = v\,\tau_z\otimes(\sigma_x p_x + \sigma_y p_y + \sigma_z p_z)
      + m\,\tau_x \otimes \bbI .
\label{eq:HDcanon}
\end{equation}

\section{Objetos dinamicos}

\(\Hop\) denota un Hamiltoniano efectivo o continuo; \(\Uop\) denota el paso
unitario discreto exacto. En particular \UW\ es la caminata Weyl minima y
\UD\ la caminata Dirac de cuatro componentes.

\section{Tres niveles}

\begin{description}
  \item[\NivelI] red discreta, localidad de \Uop, grupo de simetrias
        discretas exactas, periodicidad de cuasienergia.
  \item[\NivelII] linealizacion en un nodo, bloques Weyl/Dirac, isotropia
        aproximada, masa y velocidad efectivas.
  \item[\NivelIII] \gmu, algebra de Clifford, generadores de Lorentz,
        $SO^+(3,1)$, Poincare, $SL(2,\mathbb{C})$ como \emph{organizacion
        del limite efectivo}.
\end{description}

\section{Convenciones de lenguaje}

\begin{itemize}[leftmargin=*]
  \item \enquote{exacto}: simetrias de \Uop\ para todo \vk; identidades
        algebraicas cerradas; resultados independientes de expansion en \vq.
  \item \enquote{efectivo}: Hamiltonianos linealizados; isotropia emergente;
        Lorentz/Poincare como simetrias del limite continuo.
  \item \enquote{protoespacio}: la tupla \Proto, nunca una sola de sus
        proyecciones.
\end{itemize}
